% ============================================================================
%  AIMO Interpretability Challenge (NeurIPS 2026) -- technical report
%  Theme: negative results + evaluation methodology (directive #FINAL, 2026-11-02).
%
%  STYLE FILE: the official NeurIPS style is not vendored here (we do not commit
%  third-party files). Drop neurips_2026.sty next to this file and switch the
%  two lines marked [NEURIPS] below. Until then this compiles as an article at
%  roughly the same density.
%
%  HOUSE RULE FOR THIS FILE: no number appears here that is not in a committed
%  log under experiments/results/. The commit that produced it is cited inline
%  with \cmt{}. Placeholders are \TODO{} and must stay visible until filled.
%
%  PAGE BUDGET: <= 5 pages for Sections 1-7. The AI-usage statement,
%  Limitations, References and Appendix are excluded by the call for reports.
% ============================================================================
\documentclass[10pt]{article}
% [NEURIPS] replace the two lines below with:  \usepackage{neurips_2026}
\usepackage[margin=1in]{geometry}
\usepackage{times}

\usepackage[utf8]{inputenc}
\usepackage[T1]{fontenc}
\usepackage{booktabs}
\usepackage{amsmath,amssymb}
\usepackage{url}
\usepackage{xcolor}
% hyperref (which defines \texorpdfstring) is loaded by the NeurIPS style but not by our
% fallback preamble. \providecommand keeps this file compiling either way and changes no layout.
\providecommand{\texorpdfstring}[2]{#1}
% 표 주변 여백을 줄여 내용을 버리지 않고 자리를 번다 (counted 절이 5쪽 제한이다)
\setlength{\textfloatsep}{8pt plus 2pt minus 2pt}
\setlength{\intextsep}{8pt plus 2pt minus 2pt}
\setlength{\abovecaptionskip}{4pt}
\setlength{\belowcaptionskip}{2pt}
\providecommand{\checkmark}{\ensuremath{\surd}}
\newcommand{\TODO}[1]{\textcolor{red}{[TODO: #1]}}
% commit citation: \cmt{hash}{what it produced}
\newcommand{\cmt}[2]{{\footnotesize\textsf{[}\texttt{#1}\textsf{:} #2\textsf{]}}}
\newcommand{\X}{\checkmark}

\title{Measuring a measurement: self-generated labels overestimate\\
probe performance on reasoning robustness}
\author{Anonymous Author(s)\\\texttt{Submission under anonymous review}}
\date{}

\begin{document}
\maketitle

\begin{abstract}
We pre-registered and ran more than twenty hypotheses for predicting whether a model's correct
answer to an olympiad problem is \emph{robust} (survives answer-preserving perturbation) or
\emph{spurious}. Almost none cleared its null, and the one rule that did is unaffordable at test
time. The useful result is why the evaluation is uninformative. We generated 93 labels at $k{=}8$
samples; the organisers' labels resolve \texttt{max\_drop} to $0.01$, implying $\approx 100$. On our
labels a mid-generation hidden-state probe reaches AUC $0.872$; the same features under the same
protocol reach $0.573$--$0.610$ on theirs, by transfer and by direct training. We rule out the
classifier (five variants, learned parameters $24 \to 0$) and the label rule (our thresholds
reproduce theirs 82/82). At $n{=}33$ the shuffle-null band is $\approx 0.41$ wide and nothing clears
it on accuracy, problem length included. We add three cheap checkpoints that each caught a
self-deception: a length control that unmasked our first probe, a stability test showing
$\approx 1/4$ of our labels flip on re-run, and a three-parameter fit whose below-chance score was
nearly read as one feature's direction. Report the resolution and re-run stability of
self-generated labels next to any probe score obtained on them.
\end{abstract}

\section{Introduction}
\label{sec:intro}
The task gives a method a \texttt{model\_id} and a problem statement and asks for one bit: will this
model's correct answer survive answer-preserving perturbations? A label of \emph{spurious} marks a
problem the model answers correctly but loses under edits that preserve the answer, which makes the
target a property of the model--problem pair rather than of the text --- and is why the organisers'
own baselines read internal states (probing, $58.37\% \pm 7.5$) and sampling behaviour (CoT
uncertainty, $69.23\% \pm 10.13$) rather than surface features~\cite{aimo2026}.

Our contribution is not a better detector. We ran the obvious white-box programme --- linear probes
on hidden states before and during generation, self-consistency and memorisation signals, causal
knockouts, test-time perturbation --- under a fixed pre-registration discipline, and report it as a
negative result with a measured cause.

\paragraph{Claims.} (1) \textbf{Self-label resolution bounds what a probe score means}: our
$k{=}8$ labels resolve \texttt{max\_drop} to $0.125$, the organisers' to $0.01$; the same probe
scores $0.872$ on ours and $0.573$--$0.610$ on theirs (\S\ref{sec:resolution}). (2) \textbf{The gap
is neither the classifier nor the label rule} (\S\ref{sec:resolution}). (3) \textbf{Three cheap
checkpoints are what kept the programme honest} (\S\ref{sec:checkpoints}). (4) \textbf{At $n{=}33$
nothing beats a constant predictor on accuracy, length included} (\S\ref{sec:power}); the one
exception needs a token budget the evaluation does not grant (\S\ref{sec:acting}).

\section{Setup}
\label{sec:setup}
\paragraph{Labels.} Public training set \texttt{aimo-interp/train-main-v2}: 82 rows, 36 problems,
5 models; 71 rows carry a label (robust 52 / spurious 19) and 11 are withheld. Aggregated to problem
level by majority this gives \textbf{33 labels (robust 20 / spurious 13)}, so a constant ``robust''
predictor scores $0.606$. Labels are problem-determined: all 11 problems appearing under two or more
models carry the same label across models, so a model$\times$problem interaction is absent by
construction. We also generated 93 labels of our own (3 models, $\approx 30$ problems each) from
MATH level 5 and AIME with answer-preserving perturbations.

\paragraph{Budget.} One RTX PRO 6000 (96\,GB), \textbf{3600\,s for the entire prediction run}
(official FAQ 8; the README adds ``not one hour per problem''). This single constraint decided most
of this report: the signals we found live at generation caps of 8192 tokens and above, which a
final-set run cannot buy, and faithful replication of the label definition needs 24576 --- $187\%$
of the budget even at the smallest perturbation count.
\cmt{0fe94a2}{\texttt{ttp\_budget.txt}, \texttt{h22\_leaderboard\_budget.txt}}

\paragraph{Pre-registration protocol.} Every hypothesis was fixed before measurement in a versioned
notebook: primary metric, sign, null procedure (problem-level label shuffling, 1000 permutations,
the full pipeline re-run inside each), a length-residual companion, and the pass criterion. One
evaluation run per pre-registration; signs were never flipped after the fact, and where a result came
out opposite to the registered direction we report it opposite. Anything decided after seeing numbers
is labelled exploratory here and in the code.

\section{Self-label resolution, and what it does to a probe score}
\label{sec:resolution}
\begin{table}[t]
\centering\small
\caption{The same features and protocol, two label sources. Window-2048 mid-generation mean-pooled
hidden state at layer $L/2$, PCA-24, $\ell_2$ logistic, class-balanced, inner-CV $C$, problem-level
folds, 1000-permutation null.}
\label{tab:resolution}
\begin{tabular}{llcccc}
\toprule
Label source & \texttt{max\_drop} grid & Skywork & Olmo & DeepSeek & pooled / mean \\
\midrule
Ours, $k{=}8$ (93 rows)      & $0.125$ & $0.981$ & $0.762$ & $0.874$ & $\mathbf{0.872}$ \\
Organisers, transfer         & $0.01$  & $0.642$ & $0.627$ & $0.562$ & $\mathbf{0.610}$ \\
Organisers, direct training  & $0.01$  & $0.565$ & $0.638$ & $0.515$ & $\mathbf{0.573}$ \\
\bottomrule
\end{tabular}
\end{table}

\paragraph{The resolution gap is measured, not inferred.} The organisers' \texttt{max\_drop} takes
35 distinct values on a $0.01$ grid over $[-0.090, +1.000]$ --- six rows are negative, i.e.\ the
perturbed problem was solved \emph{more} often --- which implies on the order of 100 samples per
condition. Our $k{=}8$ grid is $0.125$: two failures out of eight in one perturbation type is the
smallest event that can cross the spurious threshold, against roughly 25 of 100 for them.
\cmt{8adae75}{\texttt{org\_label\_rule.txt}}

\paragraph{The probe score does not survive the label change.} Transfer freezes the weights trained
on our labels and applies them to the organisers' 33 problems; direct training re-fits on their
labels under leave-one-problem-out. Both land at $0.573$--$0.610$, inside their shuffle nulls (upper
bounds $0.71$ and $0.72$), against $0.872$ on our own. The problem sets do not overlap: exact-match 0
and maximum 8-gram Jaccard $0.000$, because our pool generator excludes the official problems.
\cmt{612707b}{\texttt{org\_sub2048\_transfer.txt}}, \cmt{a1e7498}{\texttt{org\_train\_report.txt}}

\paragraph{It is neither the classifier nor the label rule.} Five classifiers on identical rows,
folds and null draws span only $0.109$ (PCA-24 + logistic $0.872$ down to per-fold univariate
selection $0.640$), and reducing learned parameters makes it \emph{worse}, not better; re-drawing
rows from the same design gives a 95\% interval of width $0.49$--$0.76$, wider than any design
difference. Our reverse-engineered thresholds reproduce the organisers' \texttt{is\_robust}
\textbf{82/82} on all public rows, so the rule is right; and perturbation strength is capped by
answer preservation --- per-type collapse probability $10.9\%$ against the $40$--$50\%$ their
$27\%$ spurious rate requires, with \emph{more} perturbation types making a majority \emph{less}
likely. Details, controls and commits in the \emph{Supporting detail} section.

\section{Three checkpoints that caught us}
\label{sec:checkpoints}
Each cost one script and no GPU time, and each reversed a conclusion we had already written down.

\paragraph{(a) A length control, because the first probe was reading problem length.} Our first
probe scored cross-validated balanced accuracy $0.634$ on 137 problems and we believed it. Word
count alone scored AUC $0.75$ against the probe's $0.72$; the probe's margin correlated with length
at $\rho = +0.71$; adding the probe to length moved nothing ($0.741 \to 0.740$); and restricting to
problems of at least 15 words dropped the probe to $0.44$, at least 20 words to $0.36$. The cause
was in the labels, not the model: one perturbation type (Domain) destroyed 50 of 52 problems of
twelve words or fewer while the other four types were length-independent, so that label set largely
recorded a domain rewrite breaking a short problem.
\cmt{64c6ca3}{\texttt{length\_confound.py}, \texttt{topic\_confound.py}}

Residualising a confound is the usual response and it is weaker than it looks, because it also
removes real signal and the reader cannot tell which happened. The stronger test is a second data
set where the confound's sign is \emph{reversed}, and we had one: on the official 33 problems
shorter problems are robust (length$\to$robust AUC $0.669$), on our own 88 longer ones are
($0.372$). Memorisation scores followed the flip to within $0.017$--$0.035$ AUC ($0.704$ then
$0.389$), identifying them as a length proxy in a way residualisation cannot. That closed five
hypotheses at once. \cmt{64c6ca3}{\texttt{h23\_length\_proxy.txt}}

\paragraph{(b) A label-stability test, because about a quarter of our labels flip on re-run.} Scores
on small label sets move because the labels move. Re-labelling 20 problems with only the decoding
temperature changed, $0.6 \to 1.0$ with top-$k$ 40, flipped \textbf{7 of 20 labels} and pushed one
into the excluded band: agreement $65\%$. Independently, removing a detokenisation defect and
re-sampling the same two DeepSeek slices gave $62\%$ and $67\%$ agreement. Two different
interventions produced the same magnitude of change, which is what resampling variance looks like
rather than an effect of either intervention. A third measurement isolates it: two of our runs
happened to overlap on the same problems at the same settings, and on those the pure re-sample
agreement was $94\%$ at $k{=}8$ versus $74\%$ at $k{=}4$.
\cmt{261131c}{\texttt{temp\_sens\_report.txt}}, \cmt{0b72579}{\texttt{dsrelabel\_compare.txt}},
\cmt{4fcd838}{\texttt{flipent\_report.txt}}

The consequence is quantitative. A method evaluated on 30 labels of which a quarter are unstable
cannot be compared to another method at a resolution finer than that instability; our own
$k{=}4$-versus-$k{=}8$ comparison on the \emph{same} 30 problems moved one score from $0.721$ to
$0.571$ with nothing changed but the label rule's sample count. Reporting re-run agreement next to a
self-label score costs one extra run and tells the reader how much of the score is addressable.

\paragraph{(c) Do not read a multivariate fit's score as a single feature's direction.} We had
observed, five times under four different instruments, that \emph{robust} problems move \emph{more}
under answer-preserving edits than the pre-registered direction expected. One of those five was a
three-feature logistic fit on entropy-trajectory features that produced a leave-one-problem-out AUC
of $\mathbf{0.123}$ for one model --- far below chance --- and we were one step from treating
$0.123$ as strong evidence about the entropy feature's sign. It is not: in the same run the three
features \emph{individually} scored $0.531$, $0.515$ and $0.550$, i.e.\ nothing. The $0.123$ is what
a three-parameter fit does to 33 rows when it latches onto the wrong direction, and the strongest
regularisation in the grid was selected in 27 of 33 folds without preventing it.

We then pre-registered and ran the flip as a \emph{single}-feature test on four self-labelled arms it
had never been evaluated on. All four stayed inside their shuffle nulls ($0.721$, $0.562$, $0.337$,
$0.526$ against upper bounds $0.740$, $0.686$, $0.756$, $0.618$), and the model that had motivated
the flip came out at $0.337$ --- the \emph{original} direction, not the flipped one. The flip is not
supported, and the apparent evidence for it was an artefact of reading a fit as a feature.
\cmt{0fe94a2}{\texttt{entlens\_report.txt}}, \cmt{4fcd838}{\texttt{flipent\_report.txt}}

Before accepting any of these reversals we audited the labels themselves: five checks --- label
dtype, both sign conventions, text alignment, $\rho(\texttt{is\_robust}, \texttt{max\_drop})$ and an
end-to-end planted label --- all passed, so the reversals are not bookkeeping
(\emph{Supporting detail} section). \cmt{e0feda3}{\texttt{label\_audit.txt}}

\section{At \texorpdfstring{$n{=}33$}{n=33}, nothing beats a constant predictor}
\label{sec:power}
\begin{table}[t]
\centering\scriptsize
\caption{Every method we evaluated on the organisers' 33 problem-level labels, against the two
baselines. Accuracy is the decision rule's accuracy with thresholds fixed inside training folds
only; ``clears'' is accuracy above that method's own 1000-permutation shuffle-null 95\% upper bound.
H-FIN rows are exploratory (\S\ref{sec:acting}); R3 uses a different 19-problem label set and is not
directly comparable. Full per-hypothesis list in Appendix~\ref{app:negatives}.}
\label{tab:honest}
\begin{tabular}{llrrc}
\toprule
 & method & accuracy & AUC & clears own null \\
\midrule
\multicolumn{2}{l}{\textbf{baseline} \ constant ``all robust''} & $\mathbf{0.606}$ & $0.500$ & --- \\
\multicolumn{2}{l}{\textbf{baseline} \ problem length $-\log(\text{words})$, LOO threshold} & $\mathbf{0.636}$ & $0.671$ & no \\
\midrule
\multicolumn{5}{l}{\emph{probes on internal states}} \\
P1 & problem-text probe, model-mean (pre-registered primary) & $0.545$ & $0.616$ & no \\
P2 & mid-generation probe, window 2048, direct training & $0.556$ & $0.573$ & no \\
P3 & mid-generation probe, transfer with frozen weights & $0.505$ & $0.610$ & no \\
P4 & \;+ in-call rank 50\% cut instead of fixed threshold & $0.545$ & $0.610$ & no \\
P5 & \;+ 1:1 rank average with length & $0.626$ & $0.674$ & no \\
P5a & \;+ \emph{ablation}: one model's probe scores every problem & $0.485$--$0.545$ & $0.562$--$0.642$ & no \\
P5b & \;+ \emph{ablation}: per-model thresholds replaced by one & $0.495$ & $0.610$ & no \\
P6 & mass-mean direction, transfer & $0.586$ & $0.633$ & no \\
P7 & per-fold univariate selection (exploratory) & $0.656$ & $0.631$ & no \\
\midrule
\multicolumn{5}{l}{\emph{behaviour read from short generations}} \\
B1 & entropy trajectory, 3 features, LOPO & $0.424$ & $0.413$ & no \\
B2 & logit-lens layer of answer commitment & $0.606$ & $0.540$ & no \\
B3 & forced-answer agreement under perturbation, 512-tok CoT & $0.530$ & $0.455$ & no \\
B4 & accuracy drop vs.\ pseudo-gold, 1024-tok CoT, 3 samples & $0.535$ & $0.481$ & no \\
B5 & self-disagreement of 3 original samples & $0.475$ & $0.405$ & no \\
B6 & accuracy direction fitted on 156 held-out rows, projected once & $0.545$ & $0.590$ & no \\
B7 & accuracy drop vs.\ \textbf{gold} (19-problem set, both signs) & $0.579$ & $0.526$ & no \\
\midrule
\multicolumn{5}{l}{\emph{replicating the label definition at test time}} \\
D1 & $N{=}4$, $K{=}6$ answer-preserving perturbations, cap 24576 & $0.645$ & $0.673$ & no \\
\midrule
\multicolumn{5}{l}{\emph{natural termination} (exploratory, and see \S\ref{sec:acting})} \\
H1 & cap 12288, DeepSeek & $0.485$ & $0.575$ & no \\
H2 & cap 12288, Skywork & $0.697$ & $0.683$ & at bound \\
H3 & cap 12288, Olmo & $\mathbf{0.818}$ & $0.850$ & \X \\
H4 & \;3-model mean & $0.667$ & --- & \X \\
H5 & \;per-model routing hybrid (H2, H3 + length elsewhere) & $\mathbf{0.682}$ & --- & \X \\
\bottomrule
\end{tabular}
\end{table}

Table~\ref{tab:honest} is the whole programme on one axis, and three things read off it.
\emph{First, the constant predictor is hard to beat and most of our work does not beat it} ---
fifteen of twenty rows sit below $0.606$, and the AUC column agrees, with seven rows below $0.5$.
\emph{Second, problem length is the honest baseline and it is not significant either}: AUC $0.671$,
the highest of any non-generation method, yet inside its own shuffle null (upper bound $0.70$), with
accuracy $0.636$ one problem above the constant. Our submitted length rule therefore rests on
calibration rather than ranking --- score $-\log(\text{words})$, threshold fixed so the predicted
True rate equals the label prior ($20/33 = 0.6061$; threshold $-4.197174$, achieved rate exactly
$0.6061$) --- and it loads no model and generates no tokens, $9.15$\,ms for an 8-model, 40-problem
run, so the failure modes that voided our generation-based submissions cannot occur.
\cmt{cb9c291}{\texttt{lencal\_artifacts.txt}}
\emph{Third, the inner cross-validation keeps saying the same thing}: it picked the strongest
regularisation in the grid in 29 of 33 folds for one model and 22 of 33 for another, and switching
from our labels to the organisers' did not change that, which locates the symptom in the
dimension-to-row ratio rather than in label noise. Appendix~\ref{app:power} makes it concrete: with
a one-dimensional signal planted at known strength into the real evaluation path, our pipeline
recovers only $0.696$ from a $20\sigma$ signal at the dimensionality we use. A probe scoring $0.573$
here is inside the range where a planted signal would also have scored.
\cmt{e0feda3}{\texttt{label\_audit.txt}}

\section{Acting instead of probing: the budget wall, measured against gold}
\label{sec:acting}
If a probe cannot read the label off internal states, the remaining option is to \emph{act}: load the
judged model, perturb the problem in ways that preserve the answer, and watch the answer. We built
that as a submission and ran five pre-registered rounds (B3--B7 and D1 in Table~\ref{tab:honest}).

\paragraph{The mechanism works; the budget kills it.} Cutting the chain of thought at a cap and then
\emph{forcing} an answer makes every sequence emit an integer ($94.2$--$100\%$ extraction across four
models) for $911$\,s, $28\%$ of the budget. What the budget cannot buy is \emph{being right first}.
On 19 of our own problems where gold answers exist, the original answer at an affordable cap is
almost never correct --- accuracy $\mathbf{0.175}$, 14 of 19 at exactly zero --- while the labels for
those same problems were made at a 32768-token cap with eight samples, where accuracy is
$\mathbf{0.868}$. The gap is $\mathbf{0.693}$, and since the statistic is
$\text{base\_acc} - \min_t \text{var\_acc}_t$, a base accuracy of zero leaves nothing to drop: the
score took one distinct value across the set, so both pre-registered signs returned the identical
number. The metric did not point the wrong way; it ceased to exist.
\cmt{f79c33f}{\texttt{ttp\_r3\_report.txt}}

\paragraph{And the substitute everyone uses is quantitatively wrong.} Lacking gold answers, every
earlier test-time replication --- ours included --- scored ``did the answer change'' instead of ``did
the model get it wrong''. On those same 19 problems the substitute reports $0.842$ where gold reports
$0.175$: a $\mathbf{4.8\times}$ overestimate, because a consistently wrong model has a stable
majority answer. What a submission can buy is the stability of an error; what the label asks for is
the fragility of a correct solution. Implementing the definition faithfully ($N{=}4$, $K{=}6$
answer-preserving perturbations, majority rule, cap 24576) reaches AUC $0.673$ inside its null, and a
reduced $K{=}2$ rename-only version of itself reaches $0.675$ --- perturbation diversity moved the
AUC by $0.002$, so faithfulness was not the missing piece (\emph{Supporting detail} section).
\cmt{0fe94a2}{\texttt{h22\_official33\_skywork\_report.txt}}

\paragraph{One rule does clear its null --- and cannot be bought.} The exception is the simplest
behavioural signal in the set and not ours originally: a public competitor repository reports a
non-probing rule in which a problem is robust if its attempts terminate naturally inside the token
cap~\cite{econlearn}. We had measured that quantity as a continuous completion rate (AUC $0.671$ at
cap 16384, $0.779$ at cap 8192) but never as a decision rule, having closed the generation-length
family as budget-bound. Re-evaluated as an absolute rule --- robust iff at least half the attempts
finish inside the cap --- on stored generations at cap 12288 it gives accuracy $0.818$ for Olmo
(AUC $0.850$, null upper $0.66$), $0.697$ for Skywork and $0.485$ for DeepSeek: a three-model mean of
$\mathbf{0.667}$. Routing the two models that clear their nulls to this rule and the other two to the
length rule gives $\mathbf{0.682}$ against a null upper bound of $0.637$ --- the highest number we
measured on these labels and the only one above both baselines.
\cmt{98d20b7}{\texttt{hfin\_report.txt}, \texttt{hfin\_hybrid\_report.txt}}

Three caveats make this a measurement rather than a method. \textbf{(i) The routing is post-hoc}
--- the two models were chosen \emph{after} seeing per-model accuracies, so the pre-registrable
object is the three-model mean $0.667$. \textbf{(ii) The setting that works cannot be bought at
final-set scale}: only cap 12288 clears the gate (cap 8192 gives $0.546$, cap 4096 $0.404$), and
under our measured timing model 35 problems at that cap cost $2802$\,s for Skywork and $3650$\,s for
Olmo against a $3300$\,s total, so at final-set scale both models fall back and the submission
\emph{becomes} the length rule. Lowering the cap does not help: ``finished naturally'' saturates to
False, which is silence, not signal. \textbf{(iii) The rule may be re-reading the label definition}
--- the organisers' labels are accuracies inside a fixed budget, so ``does it finish inside a
budget'' is partly the same measurement, which breaks transfer whenever that budget differs.

The hybrid was submitted on the small track (\texttt{1026614}) and scored $\mathbf{0.57}$, equal to
the length rule and to the board median: the offline margin of $0.682$ over $0.636$ did not survive
14 items, which is what one item of resolution predicts either way.

\section{Conclusion}
\label{sec:conclusion}
Three numbers carry this report. A probe scoring $0.872$ on labels we generated scored $0.610$ when
those weights were applied to the organisers' labels and $0.573$ when refitted on them, and the
difference is traceable to resolution: our labels resolve accuracy to $1/8$, theirs to $1/100$, and
about a quarter of ours flip on re-run. The definition can be reproduced inside a submission but not
usefully, because the accuracy it needs is $0.868$ at the budget the labels were made with and
$0.175$ at the budget a submission can buy. And at $n{=}33$ our own pipeline recovers only $0.696$
from a $20\sigma$ planted signal, which bounds how strongly any null here may be read, ours included.
So: when reporting a probe score obtained on labels you generated, report the resolution of those
labels and their agreement on an independent re-run beside it, and treat the score as an upper bound
until it survives a label source you did not generate.

% ============================================================================
%  Everything below is excluded from the 5-page limit by the call for reports.
% ============================================================================
\section*{Limitations}
\label{sec:limits}
We cannot separate ``the feature does not carry this label'' from ``$n{=}33$ cannot detect it'';
that separation needs more labels, which is exactly what the resolution argument says is expensive.
Our own labels cover three of the five judged models and $\approx 30$ problems each, and they are
filtered: label generation kept only problems with $\texttt{base\_acc} \ge 0.5$, and on that subset
the accuracy axis is nearly constant (standard deviations $0.120$, $0.167$, $0.073$; for one model
every greedy continuation was correct, $32/32$). Including the problems the screening stage rejected
raises those to $0.215$, $0.205$, $0.284$. Any claim of the form ``the probe reads whether the model
knows this problem'' has to contend with a label set built by discarding the problems it does not
know. \cmt{73ca8d1}{\texttt{r5\_rows.txt}}

The mid-generation features use one greedy continuation per problem while the labels aggregate
roughly 100 stochastic ones; we tested a 4-sample average at the labels' own decoding settings and it
was \emph{worse} than greedy ($0.757$ vs $0.872$), so the asymmetry is unresolved rather than
dismissed. \cmt{58b2e7d}{\texttt{probe\_gen\_report\_samp.txt}} The model$\times$problem component is
unmeasurable on the public training set by construction, since all 11 problems appearing for more
than one model carry identical labels. The H-FIN result in \S\ref{sec:acting} is computed from stored
generations rather than a fresh pre-registered run, its routing is post-hoc, and we have not measured
it on any label source other than the organisers' 33. Two measurement defects we found late are
reported in Appendix~\ref{app:defects}; both inflated or deflated a model's apparent capability and
both were invisible in the aggregates we had been storing.

\paragraph{Why we do not lean on leaderboard position.} The small track has 14 binary items and
allows 100 submissions, so a participant can recover the entire label vector by flipping one item at
a time in about 15 submissions. At the time of writing the public board holds 1726 submissions, with
9 perfect scores and 53 above $0.86$ on the small track and 1 perfect on the main track; medians are
$0.57$ and $0.52$. We did not probe the labels and we do not report rank. Every number in this report
is measured against the public training labels or our own, never against the hidden set.
\cmt{98d20b7}{\texttt{codabench\_ids.txt}}

Finally, a reproducibility limitation of our own release: the development repository tracks 3.5\,GB
of feature caches including six files above GitHub's 100\,MB limit, so the public release is a
snapshot containing all code and all measurement logs (664 files, 13\,MB) but not the \texttt{.npz}
caches. Commit hashes cited above refer to the private development history.

\section*{AI-usage statement}
This work was carried out by the author in a continuous collaboration with Claude (Anthropic),
used as a coding and analysis agent inside a terminal agent harness. Concretely, the AI wrote
essentially all of the code in the released repository --- feature extraction, the evaluation
scripts, the label-generation pipelines, the submission implementations and the cloud job runners ---
executed the measurements, maintained the pre-registration notebook, and drafted this report
including its tables. The author set the research direction, chose and approved every hypothesis and
pre-registration before measurement, decided the budget and which experiments to buy, performed all
Codabench uploads, and made the final calls on what to claim. Verification was treated as the
author's responsibility and as an explicit standing rule for the agent: every number in this report
is printed by a committed script into a committed log, several of the corrections and retractions
recorded here originated in the author challenging an AI-stated conclusion, and the three
checkpoints in \S\ref{sec:checkpoints} were instituted after specific AI-authored misreadings were
caught. To state the boundary explicitly, the following involved no AI system: the decision of what to
pre-register, when to abandon a line of work and what to claim or retract; every Codabench
submission and every leaderboard reading; and the production of any reported number --- each
value is printed into a committed log by a committed script, and a checker confirms that every
three- and four-decimal figure in this report occurs in one of those logs (123 of 123, with
integers and two-decimal leaderboard scores outside its scope). One place where a language model
produced experimental \emph{data} rather than assistance --- the model-rewritten perturbations
in the label-generation study --- is reported as a measured step with its own acceptance test. No AI system is an author; responsibility for the content is the author's.

\begin{thebibliography}{9}
\bibitem{aimo2026}
AIMO Interpretability Challenge organisers.
\emph{Detecting spurious reasoning in mathematical problem solving}.
arXiv:2607.13899, 2026. (Task definition, label definition, and the probing and
CoT-uncertainty baselines we compare against.)

\bibitem{econlearn}
\textsc{EconLearn}. \emph{aimo-interp-difficulty-vs-fragility}. Public repository, 2026.
\url{https://github.com/EconLearn/aimo-interp-difficulty-vs-fragility}
(Non-probing natural-termination rule; reports leave-one-template-out $\approx 0.845$ on the 82
public rows and 17/21 on the main track. We reached the same quantity independently from our own
completion-rate measurement and did not use their code.)

\bibitem{byeolki}
\textsc{byeolki}. \emph{aimo-interp}. Public repository, 2026.
\url{https://github.com/byeolki/aimo-interp}
(Reports that a prompt-only probe reads solve rate rather than robustness. Our \S\ref{sec:power}
row B6 tests that reading directly and finds the solve-rate direction is not even recovered
in-domain, with cross-validated correlations $0.102$, $-0.067$, $0.172$.)
\end{thebibliography}

\appendix
\subsection*{Supporting detail (moved from the body for the page limit)}
\label{app:detail}
Verbatim measurement detail for the compressed claims in \S\ref{sec:resolution}, \S\ref{sec:checkpoints} and \S\ref{sec:acting}.

\paragraph{It is not the classifier.} If over-fitting at 30 rows were the bottleneck, reducing
learned parameters should help; it does the opposite. On identical rows, folds and null draws:
PCA-24 + logistic $0.872$, mass-mean (one direction) $0.868$, a hyperdimensional prototype with no
fitted parameters $0.824$, the same with 16-segment order encoding $0.763$, per-fold univariate
selection $0.640$. Two controls: ridge on PCA-24 and ridge on the raw $3584/4096$-dimensional
features give \emph{identical} scores (max abs.\ difference $6.0\times10^{-8}$), because at 30 rows
the feature matrix has rank $\le 29$; and the continuous \texttt{max\_drop} target instead of the
binary label lowers the score by $0.088$. The span of those five cells is $0.109$, whereas re-drawing
rows from the same design, labels and evaluation path gives a 95\% interval of width $0.49$--$0.76$
(200 draws per model) --- so we do not even claim the ordering.
\cmt{9c7a10e}{\texttt{cmp4\_report.txt}}, \cmt{f730878}{\texttt{probe\_subset\_shift.txt}}

\paragraph{It is not the label rule either.} We had reverse-engineered the organisers' thresholds as
$\texttt{max\_drop} \le 0.10 \Rightarrow$ robust, $\ge 0.25 \Rightarrow$ spurious, in between
excluded. On all 82 public rows that rule reproduces \texttt{is\_robust} \textbf{82/82}: the 11
withheld rows have \texttt{max\_drop} in $[0.110, 0.210]$, exactly the gap, and a grid sweep pins
the lower threshold at $0.10$ with only the upper loose. This is \emph{in-sample} --- the rule came
from these rows --- so it locates the thresholds, it does not validate them. The defect in our
labels is resolution and reproducibility, not the decision rule.
\cmt{8adae75}{\texttt{org\_label\_rule.txt}}

\paragraph{A cheaper symptom of the same problem.} Perturbation strength is bounded by
answer preservation: across our $k{=}8$ label sets the per-type collapse probability is $10.9\%$,
while the organisers' spurious rate of $27\%$ (19/71) under the same majority rule needs
$40$--$50\%$. Since the majority threshold rises with the type count, adding types \emph{lowers} the
chance a majority fires ($3.3\%$ at $T{=}3$, $0.2\%$ at $T{=}6$), and a model-rewritten Domain
perturbation passed its checks for only $9/20$ problems against a pre-registered $60\%$ gate.
\emph{The perturbations we can build under answer preservation do not make the organisers' rule
fire.} \cmt{6e76312}{\texttt{strong\_need.txt}}

\paragraph{We audited the labels before accepting any reversal.} Five checks all passed --- no
string values in the label column (so the classic \texttt{bool('False')} inversion cannot occur),
both sign conventions verified on hand-built scores, $33/33$ text alignment (which also exposed that
only $18/33$ problem indices coincide with label indices --- a real misalignment we had hit and
fixed), $\rho(\texttt{is\_robust}, \texttt{max\_drop}) = -0.773$, and a planted label recovered
through the real nested-CV code. \cmt{e0feda3}{\texttt{label\_audit.txt}}

\paragraph{The mechanism works; the measurement does not.} The design cuts the chain of thought at
a cap and then \emph{forces} an answer, so every sequence emits an integer however far the reasoning
got: integer extraction was $94.2$--$100\%$ across four models and three cut points, and a
four-model, forty-problem run costs $911$\,s --- $28\%$ of the budget, against the $187\%$ the same
design needs at a 2048-token cap. \cmt{0fe94a2}{\texttt{ttp\_report.txt}}

\paragraph{Why it fails: the budget sets the accuracy, and the accuracy sets the metric.} Round B7
ran on 19 of our own problems for which gold answers exist, so for once the definition could be
evaluated without substituting self-consistency for correctness. At the cap a submission can afford
the model's original answer is almost never right: accuracy $\mathbf{0.175}$, median $0.000$, 14 of
19 problems at exactly zero --- while the same problems' labels were generated at a 32768-token cap
with eight samples, where accuracy is $\mathbf{0.868}$. The gap is $\mathbf{0.693}$. Since the
statistic is $\text{base\_acc} - \min_t \text{var\_acc}_t$, a base accuracy of zero leaves nothing
to drop: the score took one distinct value across the set, which is why both pre-registered signs
produced the identical number. The metric did not point the wrong way; it ceased to exist.
\cmt{f79c33f}{\texttt{ttp\_r3\_report.txt}}

\paragraph{And the substitute everyone uses is quantitatively wrong.} Every earlier test-time
replication of the definition, ours included, replaced ``did the model get it wrong'' with ``did the
answer change'', because the public training set has no answer column. On the same 19 problems that
substitute reports $0.842$ where gold reports $0.175$: a $\mathbf{4.8\times}$ overestimate. The
mechanism is visible once stated --- a model that is \emph{consistently} wrong has a stable majority
answer, so self-consistency is high exactly where accuracy is zero. Rounds B3 and B4 were therefore
measuring the stability of a wrong answer, and B5 shows the same from the other side: in 11 of 19
problems all three original samples agree, and those answers are mostly wrong. \emph{What a
submission can buy is the stability of an error; what the label asks for is the fragility of a
correct solution.} \cmt{78709ee}{\texttt{ttp\_r4\_report.txt}}

\paragraph{Faithful replication is possible and still does not predict.} Round D1 implemented the
definition directly ($N{=}4$ samples, $K{=}6$ answer-preserving perturbations including model-written
paraphrases that pass a number-multiset check, majority rule, truncated samples excluded, cap 24576).
It reaches AUC $0.673$ inside its null $[0.31, 0.70]$. The informative comparison is against a
reduced version of itself: with $K{=}2$ and rename-only perturbations the AUC is $0.675$. Raising
diversity from two perturbation types to six moved the AUC by $0.002$, so faithful replication is
not the missing piece. \cmt{0fe94a2}{\texttt{h22\_official33\_skywork\_report.txt}}

\section{Condensed negative results}
\label{app:negatives}
\begin{table}[h]
\centering\small
\caption{All pre-registered hypotheses, one line each. ``Primary'' is the pre-registered primary
metric on the organisers' 33 problem-level labels unless noted; nulls are 1000-permutation
problem-level shuffles. ``rev.'' marks a result opposite to the registered sign, which we did not
flip. Logs are under \texttt{experiments/results/}.}
\label{tab:negatives}
\begin{tabular}{llrlc}
\toprule
 & hypothesis (primary metric) & value & null 95\% & verdict \\
\midrule
H1 & batch median split, 12-feature LOPO & --- & --- & no signal \\
H2 & difficulty $\times$ capability (label consistency) & $100\%$ & --- & labels are problem-level \\
H3 & problem pooling, shared instrument & $0.33$ & --- & length only, rev. \\
H4 & external validation, MATH$\to$v2, layer 18 & $0.73$ & $[0.36,0.75]$ & no \\
H5 & memorisation: NLL, Min-K\%, zlib & $0.69$ & --- & length proxy \\
H6 & internal shake under edits, $d_{\text{last}}$ & $0.29$ & below & rev. \\
H8 & surface-gradient dependence & $0.51$ & --- & no \\
H9 & essence separation (deletion importance) & $0.596$ & $[0.32,0.65]$ & no \\
H10 & trajectory compass, 4 models & $0.615$ & $[0.28,0.70]$ & no \\
H11/H11b & surface-token knockout, all queries & $0.403$ & --- & no (effect $\approx 0$) \\
H13 & surface-subspace PCA energy & $0.512$ & $[0.37,0.64]$ & no \\
H14/H14b & generation length ratio under rename & $0.400$ & $[0.29,0.70]$ & rev. \\
H15 & renameable surface symbol count & $0.760$ & $[0.30,0.70]$ & \X\ raw, no after length \\
H16/H16b & prefill pooled probe & $0.531$ & $[0.31,0.72]$ & no \\
H18--H19 & Min-K\%++, prefix completion, NLL gap & $0.368$--$0.490$ & --- & no \\
H20 & original-only generation length & $0.850$ & $[0.30,0.72]$ & \X\ but budget-bound \\
H20-M & per-model length, 3 models & $0.823$ & --- & saturation confound \\
H21a--d, H21p & self-uncertainty, forced-answer drop & $0.802$ & $[0.28,0.71]$ & \X\ then failed replication \\
H21cK & $K$ constant-perturbation average & $0.604$ & --- & no (single-variant noise) \\
H22 & faithful definition replication & $0.673$ & $[0.31,0.70]$ & no \\
H23 & membership inference, own-model labels & $0.408$ & $[0.37,0.63]$ & no, length proxy \\
\#77 & final probe protocol, model mean & $0.616$ & upper $0.78$ & no \\
\#GEN-PROBE & mid-generation probe, window 512 & $0.619$ & upper $0.73$ & no \\
\#CMP4 & four classifiers & span $0.109$ & --- & classifier is not the cause \\
\#LIT-PLAN & entropy trajectory, logit lens, mass-mean & $0.413$--$0.633$ & --- & no \\
\bottomrule
\end{tabular}
\end{table}

Three of these are structural rather than simple nulls, and they are the ones we would carry
forward. (i) \emph{The inference budget selects the method}: the generation-length family needs cap
$\ge 8192$, which a final-set run cannot buy, so the method that scores best offline is unavailable
online. (ii) \emph{The label definition is reproducible at test time but not predictable from
behaviour}: H22 at AUC $0.673$ inside its null, with perturbation diversity worth $0.002$.
(iii) \emph{Answer-preserving perturbations cannot make the majority rule fire}: per-type collapse
$10.9\%$ against the $40$--$50\%$ the organisers' spurious rate implies.

\section{Resolution of our own evaluation}
\label{app:power}
\begin{table}[h]
\centering\small
\caption{AUC recovered by the real probe pipeline when a signal of known size is planted in one
feature dimension, $n{=}33$ problems, PCA to 24 components, layer and $C$ chosen inside folds.}
\label{tab:power}
\begin{tabular}{rrrrr}
\toprule
feature dimension & $1\sigma$ & $3\sigma$ & $8\sigma$ & $20\sigma$ \\
\midrule
$10$ & $0.669$ & $0.996$ & $1.000$ & $1.000$ \\
$40$ & $0.715$ & $0.873$ & $0.927$ & $0.946$ \\
$400$ & $0.515$ & $0.585$ & $0.685$ & $\mathbf{0.696}$ \\
\bottomrule
\end{tabular}
\end{table}

Re-drawing rows is a second, independent read on the same ceiling. Holding the design, the labels
and the evaluation path fixed and resampling rows to the same count gives 95\% intervals of width
$0.49$--$0.76$ (200 draws per model) --- wider than any design difference we measured. That is why
\S\ref{sec:resolution} declines to rank its five classifiers.
\cmt{f730878}{\texttt{probe\_subset\_shift.txt}}

\section{Two measurement defects}
\label{app:defects}
A tokeniser-level one: vLLM detokenised one of the four models through a legacy path, leaving a
byte-level marker for spaces, so the grader scored correct answers as wrong at a rate of $0.293$.
For two weeks this read as a model difference --- lowest \texttt{base\_acc}, highest spurious rate,
below-chance AUC --- and all of it was one bug, visible only once we started saving raw generations.
We re-generated the affected labels; about one third of them changed, and the model's mean original
accuracy moved from $0.70$ to $0.97$. \cmt{abf4043}{\texttt{tok\_audit}}

An extraction-level one: the answer parser reads \texttt{ANSWER:} lines and inline
\texttt{\textbackslash(\textbackslash boxed\{\}\textbackslash)} but not display
\texttt{\textbackslash[\textbackslash boxed\{\}\textbackslash]}, costing one model $5.0\%$ of its
measured accuracy while costing another $0\%$, because the choice of output format differs by model.
\cmt{ff9cb8c}{\texttt{extract\_fail\_audit.txt}}

Both defects inflate or deflate a model's apparent capability, both were invisible in the aggregates
we had been storing, and both are the reason our standing rule is now that any generation run must
save its raw text.

\section{Submissions}
\label{app:subs}
Methods were packaged as Codabench submissions. The per-track leaderboard shows only the latest
submission per track, so the first two IDs below were read from the account's submission history
rather than from the public board.

\begin{center}\small
\begin{tabular}{llrl}
\toprule
submission\_id & track & accuracy & method \\
\midrule
\texttt{979457} & Small ($N{=}14$) & $\mathbf{0.57}$ & mid-generation probe, \textbf{nominated} (\S\ref{sec:resolution}) \\
\texttt{979466} & Main ($N{=}21$)  & $\mathbf{0.62}$ & mid-generation probe, \textbf{nominated} \\
\texttt{1000005} & Main ($N{=}21$)  & $0.48$ & prior-calibrated length rule (\S\ref{sec:power}) \\
\texttt{1000047} & Small ($N{=}14$) & $0.50$ & prior-calibrated length rule \\
\texttt{1005084} & Small ($N{=}14$) & $0.50$ & front-of-reasoning entropy, experimental \\
\texttt{1026614} & Small ($N{=}14$) & $0.57$ & natural-termination hybrid, exploratory (\S\ref{sec:acting}) \\
\bottomrule
\end{tabular}
\end{center}

The natural-termination hybrid of \S\ref{sec:acting} is the only method in
Table~\ref{tab:honest} above both baselines, and it was packaged for both tracks after the
measurements above (\texttt{dist/hfin-main.zip}, \texttt{dist/hfin-small.zip}; routing: Olmo and
Skywork use the rule with a 12288\,$\to$\,8192 per-case budget ladder, every other
\texttt{model\_id} including \texttt{gpt-oss-120b} at all effort levels takes the length rule
without loading a model). At Main scale, four to five problems per model, the first rung costs
508--582\,s per generating call against a $\approx 1650$\,s share, i.e.\ $31$--$35\%$ of the budget.
It was submitted on the small track before the deadline as \texttt{1026614} and scored $\mathbf{0.57}$ (8/14)
--- equal to the length rule's $0.57$ and to the board median, i.e.\ one item apart and not
distinguishable from it. We record it as an exploratory submission, not a nomination: the
routing that produced the offline $0.682$ was chosen after seeing per-model accuracies. We hold
no worker log for that run, so we cannot state which rung of the budget ladder executed; our
timing model admits cap 12288 at leaderboard scale (six generating problems, $57\%$ of budget),
which is the rung the offline number comes from.

\paragraph{Nominated submission.} We nominate the mid-generation probe that
Table~\ref{tab:resolution} is about (\texttt{dist/sub2048.zip}), the one white-box method we did
submit: \texttt{979466} on Main, scoring $0.62$ (13/21), and \texttt{979457} on Small, scoring
$0.57$ (8/14). Both runs were later overwritten on the per-track leaderboard, which displays only
the latest submission per track; the IDs come from the account's submission history.

Two measurements belong beside that nomination, and both are in the body rather than hidden here.
First, its fixed threshold is miscalibrated off its training distribution: the predicted True rate
is $0.48$--$0.50$ on the labels it was fitted to and $0.15$--$0.21$ on the official 33, where its
accuracy falls below the constant predictor. Second, the same frozen weights transfer at pooled AUC
$0.610$, below the $0.671$ of problem length alone. Removing either per-model partition --- one
instrument for every model, or one shared threshold --- does not recover anything (rows P5a and P5b
of Table~\ref{tab:honest}); so whatever the Main score of $0.62$ was, it was not the per-model
routing. We nominate it because it is the submission this report analyses end to end, not because
the number is defensible.

\subsubsection*{Reproducibility}
Every number cites the commit that produced it; logs live in \texttt{experiments/results/} and the
pre-registration for each hypothesis, including the ones that failed, is in the versioned notebook
file in the repository. Code, all measurement logs and every pre-registration are served
anonymously at
\begin{center}
\texttt{https://\allowbreak anonymous.4open.science/\allowbreak r/\allowbreak aimo-interp-log-6A8E}
\end{center}

\end{document}
